\originalchapter{质量损失、核心坍缩与脉冲星}
本章对应8.901第16--18和21--23次课中的高能恒星部分。质量决定可达阶段，但风、双星交换、组成与自转会改核心质量；因此观测类型与最终遗迹不能按一个总质量阈值机械对应。

\originalsection{核能停止补充后，压强支持为何失效}
大质量星依次烧较重核，中心形成铁峰富集核心。铁峰核每核子的结合能接近最大，后续聚变不再持续放热。高温光子可将重核拆开，消耗热能；电子俘获$p+e^-\to n+\nu_e$减少电子数与电子压，逃逸中微子又带走能量。若有效$\Gamma_1$软化、核心超过可支持质量，准静态序列无法维持，坍缩接近动力学时间。

达到核密度后EOS变硬、核心反弹和激波并不自动保证外层成功爆炸：激波会消耗能量于光解离、遇到吸积，并受中微子加热、多维流动和旋转影响。这里给的是物理预算链，未声称执行辐射流体超新星模拟。PS7.1从密度变化实际估引力释放，PS7.2则只求题设两物种核统计平衡，不能拿那张图当完整核素产额。

均匀$1.2\Ms$核收缩至几公里给$\Delta E\sim10^{53}$ erg，远大于通常可见光辐射的能量尺度；其分配到中微子、动能与光的比例是爆发物理问题，不能认为引力释放全变成表面光。热大爆炸BBN与恒星爆发核合成的密度、时间、弱平衡和核网络也不同。

\originalsection{风与双星会先改写核心的外边界}
光子动量率$L/c$可以通过散射或谱线作用于气体；比较辐射与重力给Eddington尺度，但线驱动风需要频率分布与速度梯度，恒定灰不透明度不能给所有风的精确质量损失。低温巨星还可能有尘埃、脉动和磁过程，Reimers是给定参数的经验律，PS6.5按原题固定$L,T_{\rm eff},R$求一个理想失质量过程，没有把它冒充实际恒星全部演化。

双星的Roche势由引力和共转离心势组成，临界等势面限制每颗星能保持的包层。溢出后质量与角动量可重新分配，保守转移从轻到重使轨道张开，重到轻使其收缩；稳定还需比较供体绝热半径响应与Roche瓣响应，PS9.2给完整$q$阈值推导。瞬时超新星失质量则改变新的轨道能，不能套用缓慢保守转移的守恒量。

光谱中的超新星分类看氢、氦及特征线，爆发机制分类则看核心坍缩或白矮星热核失控。两者相关但不一一相同。GRB可能与致密并合或坍缩发动机相联，喷流方向、相对论加速和辐射转化会影响看到的持续时间；单凭一次高能光变不能推出唯一发动机。

\originalsection{脉冲星是转动磁化中子星的几何和力矩问题}
转动与倾斜磁轴使辐射束扫过视线，观测脉冲不必表示整个星体周期性爆炸。粒子加速、电荷供给和磁层是辐射机制的重要环节；真空磁偶极只是一个可推的自转减慢模型。

转动能$E_{\rm rot}=I\omega^2/2$，若$I$固定，$I\omega\dot\omega=-L_{\rm loss}$。磁偶极损失$\propto\mu^2\omega^4\sin^2\alpha$给$\dot\omega=-k\omega^3$，因而$P\dot P$连接磁场归一。PS8.4区分赤道／极场、固定倾角、制动指数和初周期；$\tau_c=P/(2\dot P)$只有在相应模型与$P_i\ll P$下才近真实年龄。

自转过快时表面赤道质点可被离心作用抛离，Newton尺度$P_{\min}=2\pi\sqrt{R^3/(GM)}$；真实快速转动需GR、形状与EOS。PS8.3还展示为什么把慢转隆起公式代入很大扁率只得到粗估，而非精确平衡星形。

射电脉冲在等离子体传播的群速依频率，时间延迟$\propto{\rm DM}\,\nu^{-2}$；DM是电子柱密度，需电子空间模型才能换成距离。若脉冲星吸积，磁层、共转半径和质量流决定力矩，可自转加速或进入推进状态。PS9.3明确原题归一歧义；不能因幂次正确就假称所有几何系数已被证明。

\originalsection{引力波测的是轨道能量流失}
双星四极矩随轨道变化，远区TT波有两个偏振。圆轨道面向观察者可呈圆偏振，侧面可呈线偏振；传播方向投影决定它们，不能只用一个未投影矩阵解释所有方向。

辐射使轨道总能更负、间距和周期缩小。PS9.4从等质量圆轨道逐步计算四极矩、功率和$\dot P$，并说明实际Hulse--Taylor双星的偏心率为何造成圆模型与测量差别。本卷GR卷负责一般波动推导，天体卷用具体系统说明观测量、几何和近似条件。
